% GENERATED FILE. Edit canonical lessons, then run npm run generate:books.
% canonical-source-hash: fnv1a64:ac2b66e9a41bbee0
% canonical-lessons: UR-C456-zalim, UR-C456-sangdil, UR-C456-ehsan, UR-C456-karam, UR-C456-sitam

\chapter{Kind and Cruel}
\label{ch:ur-kind-and-cruel}
\hlchaptermodality{456}

\begin{chapteropening}
\textbf{By the end of this chapter:} \emph{I can say whether someone is kind or harsh.}

Complete the last lesson of chapter 456: I can say whether someone is kind or harsh.
\end{chapteropening}

\section[\texorpdfstring{\ur{ظالم} (zālim)}{ظالم (zālim)}]{\ur{ظالم} (zālim) — cruel; oppressive}

\label{lesson:UR-C456-zalim}



\begin{quote}
Before the new one: say the Urdu for slothful, then the Urdu for heedless.
\end{quote}

\subsection*{You'll want to know: \ur{ظالم}}
\textbf{\ur{ظالم}} — \emph{zālim} — \textquotedblleft{}cruel; oppressive\textquotedblright{}.

\begin{quote}
\ur{پرانا} \ur{بادشاہ} \ur{بڑا} \ur{ظالم} \ur{تھا۔}
\end{quote}

\begin{quote}
\emph{purānā bādshāh baṛā zālim thā.}
\end{quote}

\begin{quote}
The old king was very cruel.
\end{quote}

\begin{cousinweb}
\textbf{\ur{ظالم}} (\emph{zālim}) is usually traced, through Persian, to Arabic \emph{ẓālim}, \textquotedblleft{}a wrongdoer\textquotedblright{}.
\end{cousinweb}

\begin{grammarlens}[title={What you've built}]
One more word for describing what you lived and felt.
\end{grammarlens}

\subsection*{Your turn}
\begin{itemize}
\raggedright
  \item \emph{Say it:} \emph{zālim}
  \item \emph{Say it:} \emph{zālim}, once more
  \item \emph{Say it:} \emph{purānā bādshāh baṛā zālim thā.} — The old king was very cruel.
  \item \emph{Recall:} say the Urdu for slothful, then the Urdu for heedless, then say \emph{zālim} again
\end{itemize}

\begin{tcolorbox}[breakable,colback=teal!4,colframe=teal!35!black,title={Before you move on}]
What does \ur{ظالم} mean? (\textquotedblleft{}Cruel; oppressive\textquotedblright{}.) Say it once more.
\end{tcolorbox}
\section[\texorpdfstring{\ur{سنگدل} (sangdil)}{سنگدل (sangdil)}]{\ur{سنگدل} (sangdil) — hard-hearted}

\label{lesson:UR-C456-sangdil}



\begin{quote}
Before the new one: say the Urdu for heedless, then the Urdu for cruel.
\end{quote}

\subsection*{You'll want to know: \ur{سنگدل}}
\textbf{\ur{سنگدل}} — \emph{sangdil} — \textquotedblleft{}hard-hearted\textquotedblright{}.

\begin{quote}
\ur{سنگدل} \ur{آدمی} \ur{لوگوں} \ur{کا} \ur{دکھ} \ur{نہیں} \ur{سمجھتا۔}
\end{quote}

\begin{quote}
\emph{sangdil ādmī logoṅ kā dukh nahīṅ samajhtā.}
\end{quote}

\begin{quote}
A hard-hearted man does not understand people's sorrow.
\end{quote}

\begin{cousinweb}
\textbf{\ur{سنگدل}} (\emph{sangdil}) is usually traced to Persian \emph{sangdil}, \textquotedblleft{}stone-hearted\textquotedblright{}: \emph{sang}, \textquotedblleft{}stone\textquotedblright{}, and \textbf{\ur{دل}}, \textquotedblleft{}heart\textquotedblright{}.
\end{cousinweb}

\begin{grammarlens}[title={What you've built}]
One more word for describing what you lived and felt.
\end{grammarlens}

\subsection*{Your turn}
\begin{itemize}
\raggedright
  \item \emph{Say it:} \emph{sangdil}
  \item \emph{Say it:} \emph{sangdil}, once more
  \item \emph{Say it:} \emph{sangdil ādmī logoṅ kā dukh nahīṅ samajhtā.} — A hard-hearted man does not understand people's sorrow.
  \item \emph{Recall:} say the Urdu for heedless, then the Urdu for cruel, then say \emph{sangdil} again
\end{itemize}

\begin{tcolorbox}[breakable,colback=teal!4,colframe=teal!35!black,title={Before you move on}]
What does \ur{سنگدل} mean? (\textquotedblleft{}Hard-hearted\textquotedblright{}.) Say it once more.
\end{tcolorbox}
\section[\texorpdfstring{\ur{احسان} (ehsān)}{احسان (ehsān)}]{\ur{احسان} (ehsān) — a kindness done; a favour}

\label{lesson:UR-C456-ehsan}



\begin{quote}
Before the new one: say the Urdu for cruel, then the Urdu for hard-hearted.
\end{quote}

\subsection*{You'll want to know: \ur{احسان}}
\textbf{\ur{احسان}} — \emph{ehsān} — \textquotedblleft{}a kindness done; a favour\textquotedblright{}.

\begin{quote}
\ur{پڑوسی} \ur{نے} \ur{مشکل} \ur{وقت} \ur{میں} \ur{ہم} \ur{پر} \ur{احسان} \ur{کیا۔}
\end{quote}

\begin{quote}
\emph{paṛosī ne mushkil vaqt meṅ ham par ehsān kiyā.}
\end{quote}

\begin{quote}
Our neighbour did us a kindness in hard times.
\end{quote}

Masculine.

\begin{cousinweb}
\textbf{\ur{احسان}} (\emph{ehsān}) is usually traced, through Persian, to Arabic \emph{iḥsān}, \textquotedblleft{}doing good\textquotedblright{}.
\end{cousinweb}

\begin{grammarlens}[title={What you've built}]
One more word for describing what you lived and felt.
\end{grammarlens}

\subsection*{Your turn}
\begin{itemize}
\raggedright
  \item \emph{Say it:} \emph{ehsān}
  \item \emph{Say it:} \emph{ehsān}, once more
  \item \emph{Say it:} \emph{paṛosī ne mushkil vaqt meṅ ham par ehsān kiyā.} — Our neighbour did us a kindness in hard times.
  \item \emph{Recall:} say the Urdu for cruel, then the Urdu for hard-hearted, then say \emph{ehsān} again
\end{itemize}

\begin{tcolorbox}[breakable,colback=teal!4,colframe=teal!35!black,title={Before you move on}]
What does \ur{احسان} mean? (\textquotedblleft{}A kindness done; a favour\textquotedblright{}.) Say it once more.
\end{tcolorbox}
\section[\texorpdfstring{\ur{کرم} (karam)}{کرم (karam)}]{\ur{کرم} (karam) — grace; generosity}

\label{lesson:UR-C456-karam}



\begin{quote}
Before the new one: say the Urdu for hard-hearted, then the Urdu for a kindness done.
\end{quote}

\subsection*{You'll want to know: \ur{کرم}}
\textbf{\ur{کرم}} — \emph{karam} — \textquotedblleft{}grace; generosity\textquotedblright{}.

\begin{quote}
\ur{یہ} \ur{سب} \ur{آپ} \ur{کا} \ur{کرم} \ur{ہے۔}
\end{quote}

\begin{quote}
\emph{yih sab āp kā karam hai.}
\end{quote}

\begin{quote}
All this is thanks to your kindness.
\end{quote}

Masculine.

\begin{cousinweb}
\textbf{\ur{کرم}} (\emph{karam}) is usually traced, through Persian, to Arabic \emph{karam}, \textquotedblleft{}generosity, nobility\textquotedblright{}.
\end{cousinweb}

\begin{grammarlens}[title={What you've built}]
One more word for describing what you lived and felt.
\end{grammarlens}

\subsection*{Your turn}
\begin{itemize}
\raggedright
  \item \emph{Say it:} \emph{karam}
  \item \emph{Say it:} \emph{karam}, once more
  \item \emph{Say it:} \emph{yih sab āp kā karam hai.} — All this is thanks to your kindness.
  \item \emph{Recall:} say the Urdu for hard-hearted, then the Urdu for a kindness done, then say \emph{karam} again
\end{itemize}

\begin{tcolorbox}[breakable,colback=teal!4,colframe=teal!35!black,title={Before you move on}]
What does \ur{کرم} mean? (\textquotedblleft{}Grace; generosity\textquotedblright{}.) Say it once more.
\end{tcolorbox}
\section[\texorpdfstring{\ur{ستم} (sitam)}{ستم (sitam)}]{\ur{ستم} (sitam) — cruelty; injustice}

\label{lesson:UR-C456-sitam}



\begin{quote}
Before the new one: say the Urdu for a kindness done, then the Urdu for grace.
\end{quote}

\subsection*{You'll want to know: \ur{ستم}}
\textbf{\ur{ستم}} — \emph{sitam} — \textquotedblleft{}cruelty; injustice\textquotedblright{}.

\begin{quote}
\ur{اس} \ur{پر} \ur{بڑا} \ur{ستم} \ur{ہوا۔}
\end{quote}

\begin{quote}
\emph{us par baṛā sitam huā.}
\end{quote}

\begin{quote}
A great injustice was done to him.
\end{quote}

Masculine.

\begin{cousinweb}
\textbf{\ur{ستم}} (\emph{sitam}) is usually traced to Persian \emph{sitam}, \textquotedblleft{}oppression\textquotedblright{}.
\end{cousinweb}

\begin{grammarlens}[title={What you've built}]
One more word for describing what you lived and felt.
\end{grammarlens}

\subsection*{Your turn}
\begin{itemize}
\raggedright
  \item \emph{Say it:} \emph{sitam}
  \item \emph{Say it:} \emph{sitam}, once more
  \item \emph{Say it:} \emph{us par baṛā sitam huā.} — A great injustice was done to him.
  \item \emph{Recall:} say the Urdu for a kindness done, then the Urdu for grace, then say \emph{sitam} again
\end{itemize}

\begin{tcolorbox}[breakable,colback=teal!4,colframe=teal!35!black,title={Before you move on}]
What does \ur{ستم} mean? (\textquotedblleft{}Cruelty; injustice\textquotedblright{}.) Say it once more.
\end{tcolorbox}
